% chapters/00_abstract.tex -- writer R3 (W6). Abstract text only (no \chapter; main.tex supplies the heading).
Ultrathin tungsten-doped indium oxide (IWO) channels are candidates for back-end-of-line thin-film transistors, but the physical origin of their thickness dependence has not been established. This report presents a physics-constrained Silvaco ATLAS drift-diffusion model of bottom-gate IWO TFTs (TiN/\HfO/\AlO{} gate stack, Pd contacts, $L = 20~\mu$m) with channel thicknesses of 2, 6.3 and 13.2~nm, based on one measured device per thickness, and it examines which mechanisms the data can and cannot identify. The model combines an exponential acceptor tail, interface and deep states, a fixed front charge, an effective donor density, a confinement shift of the conduction band taken from first-principles slab calculations on pure \InO, and a single constant band mobility for each film. Fifty-two ATLAS launches were recorded, of which 49 produced scored currents; they include one-at-a-time parameter controls, six hypothesis tests for the 6.3~nm film and a converged recalibration (DOS 384/192 levels). The measured devices follow the pattern $2~\mathrm{nm}\approx6.3~\mathrm{nm}\neq13.2~\mathrm{nm}$, with constant-current thresholds of 0.662, 0.644 and 0.083~V and field-effect mobilities of 12.2, 11.0 and 50.2~cm$^2$/V\,s. This pattern cannot be described by a power law in thickness. On a single transfer curve, the confinement shift, the gate work function, the electron affinity and the fixed charge act together as one rigid offset, and a single shared offset fits the thresholds equally well with and without confinement (rms 0.136 and 0.133~V, respectively). Confinement is therefore an assumed input and not a result of the study. Within the calibrated trap-limited model, the approximately 4.6-fold increase in apparent mobility between 6.3 and 13.2~nm cannot be attributed to contact resistance, roughness, electrostatics or tail filling. Its origin is not determined, and a structural transition is adopted as the working hypothesis. The shared-parameter model misses the 6.3~nm threshold by $-0.294$~V, and no tested variant reproduces the on-state shape of this device without degrading the subthreshold slope, which points to a limitation of the constant-mobility model form. The model is calibrated, not validated. Output characteristics and transfer curves at 338 and 358~K were registered with SHA-256 hashes before any such data were available. They were registered in two exact variants, because ATLAS had applied a default temperature exponent for the mobility, and they define explicit falsification rules. The report closes with a list of the measurements that are required before claims about mechanisms can be made.
